\section{Convergent log-height expansion and a controlled inverse}
\label{sec:expansion}
This section incorporates a separately completed and independently
audited extension \cite{expansion,expansionaudit}. The forward result
is an absolutely convergent identity in five dependent small arguments.
The inverse result combines a convergent algebraic cubic inverse with
the first exponentially small correction. It is not a complete
all-orders exponential inverse expansion.

\subsection{An exact identity and its convergent expansion}
Keep the displayed completion, and write
\begin{gather*}
 \lambda_A=A+\sqrt{A^2-1},\quad
 c=\frac{\lambda_A^p-\lambda_A^{-p}}{2\sqrt{A^2-1}},\quad
 S=\frac{(\lambda_A^p-\lambda_A^{-p})^2}{4},\\
 \lambda_S=S+\sqrt{S^2-1},\quad
 B_*=2p(R-1),\quad a_*=A^{-2},\quad b_*=S^{-2},\quad
 u_*=\lambda_A^{-2p},\quad v_*=\lambda_S^{-2R}.
\end{gather*}
The symbols $B_*,a_*,b_*$ are local analytic variables, not the
compiler radix or fixed compiler ports. Let
\[
 g(z)=\ln\frac{1+\sqrt{1-z}}2\qquad(0\le z<1),\qquad
 c_j=\frac{\binom{2j}{j}}{2j\,4^j}.
\]
All of $d_0,a_*,b_*,u_*,v_*$ belong to $(0,1)$.

\begin{theorem}[Exact forward log-height formula]\label{thm:forward}
The exact identity
\begin{align}
 (\ln2)(\log_2\mathcal H-N)
 &=B_*\bigl[\ln(1+d_0)+g(a_*)\bigr]
   +2(R-1)\ln(1-u_*)\nonumber\\
 &\quad+R g(b_*)-\tfrac12\ln(1-b_*)+\ln(1-v_*)\label{eq:exactlog}
\end{align}
has the absolutely convergent all-orders expansion
\begin{align}
 (\ln2)(\log_2\mathcal H-N)
 =\sum_{j\ge1}\biggl[&B_*\frac{(-1)^{j+1}d_0^j}{j}
 -B_*c_j a_*^j-\frac{2(R-1)u_*^j}{j}\nonumber\\
 &+\Bigl(\frac1{2j}-Rc_j\Bigr)b_*^j-\frac{v_*^j}{j}\biggr].
 \label{eq:series}
\end{align}
After truncation at $j=k\ge1$, the absolute error in natural-log
units is at most
\begin{align}
 \mathcal R_k={}&\frac{B_*d_0^{k+1}}{k+1}
 +\frac{B_*a_*^{k+1}}{2(k+1)(1-a_*)}
 +\frac{2(R-1)u_*^{k+1}}{(k+1)(1-u_*)}\nonumber\\
 &+\frac{(R+1)b_*^{k+1}}{2(k+1)(1-b_*)}
 +\frac{v_*^{k+1}}{(k+1)(1-v_*)}.\label{eq:tail}
\end{align}
Divide this bound by $\ln2$ for the error in $\log_2\mathcal H$.
\end{theorem}
\begin{proof}
The identities
\[
 \lambda_A=2^L(1+d_0)e^{g(a_*)},\qquad
 2S=\lambda_A^{2p}(1-u_*)^2/2,
\]
and the auxiliary Pell formula give
\[
 \ln\mathcal H=(R-1)\ln(2S)+R g(b_*)
 -\tfrac12\ln(1-b_*)+\ln(1-v_*).
\]
Substitution proves \eqref{eq:exactlog}, including the coefficient
$R$ on $g(b_*)$ and the denominator correction.
Differentiating gives
\[
 g'(z)=-\frac{(1-z)^{-1/2}-1}{2z},
\]
so integration of the convergent binomial series with $g(0)=0$
yields $g(z)=-\sum_{j\ge1}c_jz^j$. The logarithm series then
prove \eqref{eq:series}; each component is absolutely convergent.
For \eqref{eq:tail}, use the alternating remainder for $\ln(1+d_0)$,
$0<c_j\le1/(2j)$, and geometric tail bounds for the other arguments.
For the $b_*$ coefficient the triangle inequality gives
$|1/(2j)-Rc_j|\le(R+1)/(2j)$, requiring no guessed sign.
\end{proof}

This formula separates main and auxiliary radical corrections from
their conjugate-power corrections. Grouping them by $j$ does not
assert a single global asymptotic ordering in the ordinary input, nor
remove the radix jumps. The arguments are dependent; that dependence
does not affect convergence of the finite sum of convergent series.

\begin{proposition}[Uniform leading correction]\label{prop:leading}
For the family, and for the real interpolation below,
\begin{equation}\label{eq:leading}
 \left|\log_2\mathcal H-N-\frac{B_*d_0}{\ln2}\right|
 <\frac{B_*d_0^2}{\ln2}.
\end{equation}
In particular $\log_2\mathcal H-N>0$, and
\[
 \log_2\mathcal H-N=
 \frac{2p(R-1)}{\ln2}\,2^{-p}
 \bigl(1+2^{1-y}+O(d_0)\bigr).
\]
The constant-one absolute remainder is the explicit statement
\eqref{eq:leading}, before this rearrangement.
\end{proposition}
\begin{proof}
Under $p\ge16$, $y\ge3$, $R\ge8$, $R<2p$, one has
\begin{equation}\label{eq:smallarguments}
 a_*,u_*,b_*,v_*\le d_0^2/16.
\end{equation}
Indeed $a_*<2^{-2p-2y}$ and $d_0\ge2^{-p}$.
Also $\lambda_A>2XY+1=2^L+1$, so
$u_*<2^{-2pL}$. For real as well as integral $p$, convexity gives
\[
 \lambda_A^{2p}>2^{2pL}+2p\,2^{(2p-1)L}>2^{2pL}+2,
\]
whence $S>2^{2pL-2}$, $b_*<2^{-4pL+4}$, and
$v_*<2^{-2R(2pL-2)}$ since $\lambda_S>S$.
Each of the latter three is at most $2^{-2p-4}$ by the stated
hypotheses, proving \eqref{eq:smallarguments}.

Use $|\ln(1+d_0)-d_0|\le d_0^2/2$,
$|g(z)|\le z/[2(1-z)]$, and
$|\ln(1-z)|\le z/(1-z)$. All four small arguments are below
$1/2$. In \eqref{eq:exactlog}, the five error contributions after
division by $B_*d_0^2$ are bounded respectively by
\[
 \frac12,\qquad\frac1{16},\qquad\frac1{8p},\qquad
 \frac1{16p},\qquad\frac18.
\]
Their sum is at most $179/256<1$ for $p\ge16$.
This proves \eqref{eq:leading}; positivity follows from $d_0<1$.
\end{proof}

\subsection{Real analytic interpolation and the cubic branch}
For each branch separately, extend $p,y$ to affine real functions on
$R\ge100$:
\[
\begin{array}{c|cc}
&p(R)&y(R)\\\hline
\mathrm A+&2(R+1)/3&(R+1)/3-1\\
\mathrm A-&2(R-1)/3&(R-1)/3-1\\
\mathrm B&5(R+1)/7&15(R+1)/28-1
\end{array}
\]
Use the positive real formulas above to define $S(R)$ and
\[
 \mathcal H(R)=\frac{\sinh(R\theta(R))}{\sinh\theta(R)},\qquad
 \theta(R)=\operatorname{arcosh}S(R)>0.
\]
At the integer family parameters this agrees exactly with the Pell
witness. It is an analytic interpolation, not a claim that noninteger
parameters are chart witnesses.

$A,\lambda_A,\lambda_A^p$, and then $S$ strictly increase with $R$.
For fixed $\theta>0$, the hyperbolic quotient strictly increases
with $R>1$. For fixed $R>1$, its logarithmic derivative in $\theta$
is $R\coth(R\theta)-\coth\theta>0$, since
\[
 \frac{d}{dz}(z\coth z)
 =\frac{\sinh z\cosh z-z}{\sinh^2z}>0\qquad(z>0).
\]
Thus $\mathcal H(R)$ and $\log_2\mathcal H(R)$ are strictly
increasing on each branch, independently of integer bit-length results.

Let $N(R)$ be the cubic bit length minus one:
\begin{equation}\label{eq:Ncubics}
\begin{array}{ll}
\mathrm A+:&(4R^3+4R^2-7R-1)/3,\\
\mathrm A-:&(4R^3-12R^2+9R-1)/3,\\
\mathrm B:&(25R^3+25R^2-39R-11)/14.
\end{array}
\end{equation}
Direct differentiation gives, uniformly for $R\ge100$,
\begin{equation}\label{eq:derivativebudget}
 N'(R)>3R^2,\quad 0<N''(R)<12R,\quad
 B_*(R)<2R^2,\quad 0<B_*'(R)/B_*(R)\le3/R.
\end{equation}
For example $B_*$ is respectively
$\frac43(R^2-1)$, $\frac43(R-1)^2$, and $\frac{10}{7}(R^2-1)$.
The cubic is increasing and unbounded on this interval.

\begin{theorem}[One exponential inverse correction]\label{thm:inverse}
Fix an actual or interpolated $R\ge100$, write
$\ell_h=\log_2\mathcal H(R)$, and let $r_0\ge100$ be the
unique root $N(r_0)=\ell_h$. Put $p_0=p(r_0)$, $y_0=y(r_0)$,
and
\[
 \mathcal T(r)=\frac{2p(r)(r-1)}{(\ln2)N'(r)}\,2^{-p(r)}.
\]
Then
\begin{equation}\label{eq:inverse64}
 R=r_0-\mathcal T(r_0)+\varepsilon_h,\qquad
 |\varepsilon_h|\le64\bigl(2^{-p_0-y_0}+2^{-2p_0}\bigr).
\end{equation}
\end{theorem}
\begin{proof}
Write $\ell_h=N(R)+E(R)$, where \eqref{eq:leading} makes
$E(R)>0$. Thus $r_0>R$. For $\delta_r=r_0-R$ the mean-value
theorem gives $\delta_r=E(R)/N'(\xi)$, $R<\xi<r_0$.
Using \eqref{eq:derivativebudget}, $\ln2>2/3$, and
\eqref{eq:leading} gives $0<\delta_r<2d_0<1$.
Put $q_*=2^{-p(R)}$. Since $d_0=q_*(1+2/Y)<2q_*$,
\[
 \left|E(R)-\frac{B_*(R)q_*}{\ln2}\right|
 \le\frac{B_*(R)}{\ln2}\bigl(2q_*/Y+4q_*^2\bigr).
\]
Also $B_*/((\ln2)N')<1$. Division by $N'(\xi)$ therefore
contributes error at most $2q_*/Y+4q_*^2$.
Replacing its principal denominator $N'(\xi)$ by $N'(R)$
contributes at most
\[
 \frac{B_*q_*}{\ln2}
 \frac{N'(\xi)-N'(R)}{N'(\xi)N'(R)}
 <\frac{4q_*\delta_r}{R}<\frac{20q_*^2}{R}\le q_*^2/5.
\]
Here $N'(\xi)-N'(R)<12\xi(\xi-R)$ and
$\delta_r<4q_*$ suffice; the displayed $20$ is deliberately loose.
Furthermore
\[
 0<\mathcal T(r)<2^{-p(r)},\qquad
 |\mathcal T'/\mathcal T|
 \le B_*'/B_*+N''/N'+p'\ln2<7/r+1<2.
\]
Thus changing $\mathcal T(R)$ to $\mathcal T(r_0)$ costs less
than $2q_*\delta_r<8q_*^2$. Summing the errors gives less than
$2q_*/Y+13q_*^2$.
Finally $p',y'<1$ and $\delta_r<1$ imply
$q_*<2\,2^{-p_0}$ and $Y^{-1}<2\,2^{-y_0}$.
The total is bounded by
$8\,2^{-p_0-y_0}+52\,2^{-2p_0}$, proving the stated looser
constant $64$.
\end{proof}

The full rational prefactor in $\mathcal T(r)$ tends to
$1/(3\ln2)$ in either A branch and $4/(15\ln2)$ in B.
Replacing it by this limit incurs an extra
$O(2^{-p_0}/r_0)$ term, generally too large to preserve
\eqref{eq:inverse64}. Its evaluation point and full coefficient matter.

\subsection{Convergent algebraic inverse to all orders}
Let $\gamma=4/3$ in A and $\gamma=25/14$ in B, and put
$s_h=(\ell_h/\gamma)^{1/3}$. The large cubic root has
\begin{align*}
 \mathrm A+:\quad r_0&=s_h-\tfrac13+\tfrac{25}{36}s_h^{-1}
                         -\tfrac{11}{81}s_h^{-2}+O(s_h^{-3}),\\
 \mathrm A-:\quad r_0&=s_h+1+\tfrac14s_h^{-1}
                         +0\,s_h^{-2}+O(s_h^{-3}),\\
 \mathrm B:\quad r_0&=s_h-\tfrac13+\tfrac{142}{225}s_h^{-1}
                         -\tfrac{104}{2025}s_h^{-2}+O(s_h^{-3}).
\end{align*}
Direct substitution into \eqref{eq:Ncubics} proves the displayed
coefficients. For convergence at all algebraic orders write
$N(r)=\gamma r^3+b_2r^2+b_1r+b_0$, $r=s_h h$, $w_h=1/s_h$.
The equation becomes
\[
 h^3+(b_2/\gamma)w_hh^2+(b_1/\gamma)w_h^2h
 +(b_0/\gamma)w_h^3=1.
\]
At $(w_h,h)=(0,1)$ its derivative in $h$ is three. The analytic
implicit-function theorem gives a unique convergent power series
$h(w_h)$ near zero; its recursively determined coefficients supply
all algebraic orders of the large real cubic root. No explicit
convergence radius is asserted. Theorem~\ref{thm:inverse} adds the
first controlled exponential correction beyond this algebraic branch.

